\documentclass[12pt,a4paper]{article}
\usepackage{amsmath}
\usepackage{amssymb}
\usepackage{braket}
\usepackage{listings}
\usepackage{xcolor}
\usepackage[margin=2.5cm]{geometry}

\lstdefinestyle{pseudocode}{
  basicstyle=\ttfamily\small,
  keywordstyle=\color{red!70!black},
  commentstyle=\itshape\color{red!70!black},
  morecomment=[l]{\#},
  frame=single,
  backgroundcolor=\color{gray!8},
  xleftmargin=1.5em,
  xrightmargin=1.5em,
  columns=flexible,
  keepspaces=true,
}

\title{\textbf{Foundations of Quantum Computing}\\[0.5em]
       \large Portfolio Problems}
\author{}
\date{}

\begin{document}
\maketitle
\thispagestyle{empty}

% ---------------------------------------------------------------
\section*{Week 2 --- Portfolio Problem \hfill [10 marks]}

Consider the operator
\[
  H = \begin{pmatrix} 1/\sqrt{2} & 1/\sqrt{2} \\ 1/\sqrt{2} & -1/\sqrt{2} \end{pmatrix}.
\]

\medskip
\textbf{a.}
Provide a simple \textbf{algebraic} argument that $H$ is a \textbf{reflection} when acting on 2D
vectors with real-valued coefficients.
\hfill\textbf{[2 marks]}

\smallskip
\textit{Hint}: consider some appropriate matrix product, or the effect of $H$ on certain
single-qubit states.

\medskip
\textbf{b.}
Provide a simple \textbf{geometric} argument, referring to angles between lines in a 2D diagram,
to determine the angle made by the axis of reflection, relative to the $\ket{0}$ axis.
\hfill\textbf{[2 marks]}

\smallskip
\textit{Hint}: what angle does $H\ket{0}$ make with $\ket{0}$?

\medskip
\textbf{c.}
Use the above to identify two orthogonal unit eigenvectors of $H$, and use this to write an
eigendecomposition of $H$ in ket-notation (i.e., in the form described in the Study Notes).
\hfill\textbf{[6 marks]}

% ---------------------------------------------------------------
\section*{Week 3 --- Portfolio Problems \hfill [15 marks]}

Consider the operator
\[
  \mathrm{SWAP} =
  \begin{pmatrix}1&0&0&0\\0&0&1&0\\0&1&0&0\\0&0&0&1\end{pmatrix}
\]
discussed in the Week 3 Study materials.

\medskip
\textbf{a.}
What is the inverse operation for $\mathrm{SWAP}$? From this, explain how you can infer that it
has eigenvalues $\pm 1$.
\hfill\textbf{[3 marks]}

\smallskip
\textit{Hint}: what it does to product states (e.g., computational basis states)?

\medskip
\textbf{b.}
Suppose that $\ket{\psi} = u_{00}\ket{00}+u_{01}\ket{01}+u_{10}\ket{10}+u_{11}\ket{11}$ is a
$(-1)$-eigenvector of $\mathrm{SWAP}$. Express $\ket{\psi}$ explicitly as a linear combination of
computational basis states, in ket notation and with at least one positive real-valued coefficient.
Give a justification for the coefficients presented.
\hfill\textbf{[3 marks]}

\smallskip
\textit{Hint}: what constraints does being a $(-1)$-eigenvector of $\mathrm{SWAP}$ impose on the
coefficients?

\medskip
\textbf{c.}
Because $\mathrm{SWAP}$ has eigenvalues $\pm 1$, it is also Hermitian --- and can be treated as a
measurement observable. If you were to perform a $\mathrm{SWAP}$-measurement on the state
$\ket{01}$, what would be the probability of the outcome $-1$? What is the post-measurement state
in this case?
\hfill\textbf{[6 marks]}

\medskip
\textbf{d.}
Simply write a matrix for a three-qubit coherently controlled operation, $\mathrm{CSWAP}$, in
analogy to the operators $CU$ on two qubits. Then, write a more-condensed algebraic representation
for $\mathrm{CSWAP}$, using projectors and ket-notation.
\hfill\textbf{[3 marks]}

% ---------------------------------------------------------------
\section*{Week 4 --- Portfolio Problems \hfill [25 marks]}

Suppose that Alice and Bob share two qubits $a$ and $b$ in the Bell state $\ket{\Phi^+}$. Suppose
also that Alice has a qubit $\mathsf{p}$ and Bob has a qubit $\mathsf{q}$, which are in a joint
state
\[
  \ket{\psi} = u_{00}\ket{0,0}+u_{01}\ket{0,1}+u_{10}\ket{1,0}+u_{11}\ket{1,1}\in\mathbb{C}^4
\]
that Alice and Bob don't precisely know --- and which may even be entangled.

\medskip
\textbf{a.}
Simply write, in ket notation, the state-vector for the joint state of $\mathsf{p}$, $\mathsf{a}$,
$\mathsf{b}$, $\mathsf{q}$ --- and with the qubits in \textit{that specific order}.
\hfill\textbf{[5 marks]}

\medskip
\textbf{b.}
Suppose that Alice and Bob perform the \textbf{following} operations on their qubits, initially in
the state $\ket{\psi}$:

\begin{lstlisting}[style=pseudocode]
# Alice and Bob each prepare some classical memory
# to store measurement outcomes
bit r       # for Alice
bit s       # for Bob

# What Alice does
CX p, a
Mz a -> r

# What Bob does
CX b, q
Mx b -> s
\end{lstlisting}

Describe the sequence of states produced by performing these operations, in the given order, on an
initial state of $\ket{\psi}$. In your answer, describe all of the possible outcomes of the
measurements performed, grouped together for each measurement (i.e., presenting all possible
outcomes together in your analysis).
\hfill\textbf{[15 marks]}

\medskip
\textbf{c.}
Compute the state $\ket{\psi'} = \mathrm{CX}\ket{\psi}$. What information must Alice and Bob send
each other, and what operations must they do to the qubits $\mathsf{p}$, $\mathsf{q}$, to
transform the possible states they could obtain in part~(b) to instead get $\ket{\psi'}$ for all
outcomes?
\hfill\textbf{[2 marks]}

\medskip
\textbf{d.}
Summarise, in your own words, what the significance is of the fact that Alice and Bob could manage
to compute the state $\mathrm{CX}\ket{\psi}$ in this way.
\hfill\textbf{[3 marks]}

% ---------------------------------------------------------------
\section*{Week 5 --- Portfolio Problem \hfill [25 marks]}

Let $f:\{0,1\}^2\to\{0,1\}$ be a function such that $f(x)=1$ if and only if $x=01$, and $f(x)=0$
otherwise. Let $\mathtt{O\_f}$ be a unitary oracle on three qubits, computing
\[
  \mathtt{O\_f}\ket{x_1 x_2,\,t} = \ket{x_1 x_2,\;t\oplus f(x_1 x_2)}.
\]

Motivated by Grover's algorithm (as described on page 10 of the Week 5 Study notes), consider the
following procedure:

\begin{lstlisting}[style=pseudocode]
register a: 2 bits      # a classical output register
register q: 2 qubits    # a two-qubit register
qubit t                 # an auxiliary 'target' qubit

MxReset q               # initialise q to |+,+>
MxReset t               # initialise
Z t                     # t to |->

O_f q[0], q[1], t       # oracle query
R_u q                   # reflect q about the |+,+> state

Mz q -> a               # measure the output
\end{lstlisting}

where we define the subroutine \texttt{R\_u} by the following operations:

\begin{lstlisting}[style=pseudocode]
define R_u q
    # Reflect a two-qubit state about the |+,+> state
    operand register q: 2 qubits

    H q             # Change of basis
    CZ q[0], q[1]   # Apply phase to two-qubit state
    H q             # Undo change of basis
    X q             # Interchange |0> <-> |1>
    stop
\end{lstlisting}

Using ket notation, compute the effects of the procedure above, describing the effect of the
operations in sequence. Describe the significance of the output in your own words.
\hfill\textbf{[25 marks]}

% ---------------------------------------------------------------
\section*{Week 6 --- Portfolio Problem \hfill [25 marks]}

Factorise $N=323$, using the reduction from integer factorisation to order finding, and the
reduction from order finding to phase estimation --- but using the following data to avoid having
to access a quantum computer.

In the following, take $a=2$ as the number to find the order of (rather than selecting $a$ at
random). Suppose that for this value of $a$, you performed phase estimation with the unitary
operator $U_a$, obtaining the following two samples of values of $x$ (representing eigenvalues
$e^{2\pi ix}$ of the unitary $U_a$):
\[
  x_1 = 7/24, \qquad x_2 = 5/18.
\]

Perform the following steps, showing your workings:

\medskip
\textbf{a.}
From the samples $x_1$ and $x_2$, compute a candidate for the order $1<s<N$, of $a$ modulo $N$,
describing the way in which you obtain it.
\hfill\textbf{[6 marks]}

\medskip
\textbf{b.}
What properties does $s$ have to have, to be of use in factorising $N$?
\hfill\textbf{[3 marks]}

\medskip
\textbf{c.}
From this value of $s$, find an integer $1<u<N$ such that $u^2\equiv 1\pmod{N}$, describing how
you obtain it.
\hfill\textbf{[10 marks]}

\medskip
\textbf{d.}
From this value of $u$, compute the factors of $N$, describing how you obtain them.
\hfill\textbf{[6 marks]}

\end{document}
